\chapter{Intégration dans la chaîne d'optimisation}\label{ch:integration}
\iterbanner{V1}

\section{Ce qui circule entre les phases}
\begin{figure}[htbp]
\centering
\begin{tikzpicture}[node distance=0.5cm and 0.9cm]
\node[flow] (x) {profil $(i_m,i_p,i_t)$\\ 9 bits};
\node[flow,right=of x] (aero) {\textbf{Aéro} (phase 1)\\ \code{evaluate}};
\node[flow,right=of aero,fill=cblue!6] (ser) {\textbf{SER} (phase 2)\\ \code{evaluate\_ser}};
\node[flow,right=of ser] (obj) {objectifs\\ $(C_D,\ \mathrm{SER})$\\ + contraintes};
\draw[arr] (x) -- (aero);
\draw[arr] (aero) -- node[above,font=\scriptsize] {$\alpha$} (ser);
\draw[arr] (ser) -- (obj);
\draw[arr] (aero.south) |- ++(0,-0.8) -| node[pos=0.25,below,font=\scriptsize] {$C_D$, $C_L$, validité} (obj.south);
\node[flow,above=1.1cm of ser,xshift=1.6cm,fill=corange!8] (geo) {géométrie : mêmes points NACA\\ (ou STL de \code{geometry\_export})};
\draw[arr,corange] (geo) -- (ser);
\node[flow,above=1.1cm of aero,xshift=-1.6cm,fill=cgreen!8] (sc) {scénario : $\Re$, $\alpha$ ou $C_L$,\\ $f$, corde, secteur};
\draw[arr,cgreen!60!black] (sc) -- (aero);
\draw[arr,cgreen!60!black] (sc.east) -- (geo.west);
\end{tikzpicture}
\caption{La SER dans la chaîne. Elle reçoit la géométrie et l'attitude $\alpha$ de l'aérodynamique, et un scénario radar partagé. Les optimiseurs ne voient que les deux objectifs et les contraintes.}\label{fig:flux}
\end{figure}

\paragraph{Le couplage par l'attitude.} À incidence imposée (scénario de la note d'intention), tous les profils volent à $2^\circ$ et la SER ne dépend que de la forme. À portance imposée, chaque profil a sa propre incidence (de $-5{,}5^\circ$ à $+3{,}5^\circ$ sur la grille) : l'attitude vue par le radar dépend alors de l'aérodynamique, et c'est un vrai couplage multiphysique. Le code gère les deux, car la SER reçoit toujours l'$\alpha$ calculé par l'aéro.

\section{Le code d'intégration}
\begin{codebox}[ : \code{optim/problem.py}]
\begin{lstlisting}
def simulate(idx, space, sc):
    m, p, t = space.values(idx)
    a = evaluate(m, p, t, alpha=fc.alpha, Re=fc.Re)      # ou evaluate_at_cl en mode portance
    s = evaluate_ser(m, p, t, alpha=a.alpha - offset, case=rc)   # attitude venue de l'aero
    g = [fc.t_min - t, fc.CL_min - a.CL, 0.0 if a.valid else 1.0]  # contraintes g <= 0
    return dict(..., CD_counts=a.CD * 1e4, SER_db=s.ser_db, feasible=all(gi <= 0 for gi in g))
\end{lstlisting}
\code{offset} recentre le secteur d'élévations pour un secteur asymétrique comme S2 : décaler toutes les élévations de $\varepsilon_0$ revient à remplacer $\alpha$ par $\alpha-\varepsilon_0$ dans $\shat$. \code{build\_table} évalue une fois les 512 profils et les enregistre dans \code{resultats/base\_complete\_<scénario>.csv}.
\end{codebox}

\section{Le contrat avec le groupe SER}
Pour que le vrai solveur remplace ce modèle sans rien changer ailleurs, il doit respecter l'interface suivante :
\begin{center}
\renewcommand{\arraystretch}{1.2}
\begin{tabularx}{\linewidth}{lX}
\toprule
Élément & Convention \\
\midrule
Entrée géométrique & STL de \code{geometry\_export} (mêmes points que l'aéro), repère corps : $x$ corde vers l'aval, $y$ extrados, $z$ envergure, origine au bord d'attaque \\
Entrée d'attitude & $\alpha$ (degrés), lu dans le JSON ; la direction radar est tournée, pas la géométrie \\
Scénario & fréquence, corde physique, polarisation, secteur d'élévations, à figer ensemble \\
Sortie & une valeur scalaire par profil : moyenne linéaire sur le secteur, en dB (dBsm en 3D, dBm en 2D) \\
Signature Python & \code{evaluate\_ser(m, p, t, alpha, case) -> objet avec .ser\_db} \\
Coût & compté par le même compteur que l'aéro : une « simulation » = aéro + SER \\
\bottomrule
\end{tabularx}
\end{center}

\paragraph{Procédure de remplacement.}
\begin{enumerate}
\item Écrire un adaptateur \code{evaluate\_ser} qui exporte le STL, appelle le solveur et renvoie \code{ser\_db}.
\item Comparer au modèle PO sur une dizaine de profils : on attend un accord de tendance (rôle de $t$), pas forcément de niveau.
\item Relancer \code{build\_table} pour les deux scénarios : tous les optimiseurs et toutes les figures suivent automatiquement.
\end{enumerate}

\begin{piege}
En 3D, la SER d'un tronçon d'aile extrudé contient l'écho des \emph{flancs} (les deux extrémités du tronçon), qui n'a rien à voir avec le profil. Il faut soit les retirer (\code{extrude(..., caps=False)}), soit les ignorer en ne regardant que les directions perpendiculaires à l'envergure, où ils sont vus en incidence rasante.
\end{piege}

\section{Ce que la SER impose à l'optimisation}
La SER moyenne sur le secteur est une fonction lisse de $(m,p,t)$ dans notre modèle : un modèle quadratique la représente bien, mieux que la traînée. Avec un vrai solveur, deux sources d'irrégularité apparaîtront : le passage d'une facette de l'ombre à la lumière, et les interférences entre contributions, sensibles à la fréquence. C'est l'argument de la note d'intention pour NSGA-II, qui n'utilise pas de gradient. Pour la boucle QUBO, une SER irrégulière dégraderait le modèle quadratique : il faudra le vérifier dès que les vraies valeurs seront disponibles.

\exercices
\begin{exercice}
En mode portance imposée, un profil très cambré vole à $\alpha=-4^\circ$. Dans le scénario S1, le radar le voit-il plutôt par l'extrados ou par l'intrados, par rapport à un profil qui vole à $+2^\circ$ ? Quel effet peut-on attendre sur la SER ?
\end{exercice}

\begin{pistes}
\begin{itemize}
\item Écrire dès maintenant l'adaptateur pour le format d'entrée du solveur choisi par le groupe SER, et le tester avec le modèle PO en guise de solveur factice.
\item Mettre en cache les résultats SER par profil et par scénario (comme \code{AeroSimulator}), pour que les appels au vrai solveur, plus coûteux, ne soient jamais répétés.
\item Prévoir un mode « multi-fidélité » : PO pour explorer, solveur complet pour confirmer les profils du front.
\end{itemize}
\end{pistes}
